\documentclass[12pt]{article}
\usepackage[utf8]{inputenc}
\usepackage[T1]{fontenc}
\usepackage{hyperref}
\title{Yarluyandi Language Resources for K-12 Curriculum Development}
\author{Compiled by Claude Sonnet 5.5 (claude-sonnet-5-5) for Andr\'as Kornai}
\date{}
\begin{document}
\maketitle

\section{Introduction}
Yarluyandi (ISO 639-3 yry) is a dormant Australian language of the Lake Eyre region of South Australia and Queensland. It is closely related to Wangkangurru, Ngamini and Diyari. Sources disagree on whether it is a separate language or a dialect of Ngamini. It is not taught in schools as far as I could find.

\section{Primary Educational Resources}
The Mobile Language Team portal has a Yarluyandi learning course with word lists. It was created by Anpanuwa Joyce Crombie and Aulpunda Jean Barr-Crombie, and the Wangkangurru/Yarluyandi people retain intellectual property rights over the materials, so any reuse needs their permission.
\begin{itemize}
\item \url{https://portal.mobilelanguageteam.com.au/languages/yarluyandi/}
\end{itemize}

\section{Linguistic and Reference Materials}
\begin{itemize}
\item Hercus, Luise. Yarluyandi vocabulary (unpublished wordlist, catalogued at AIATSIS as ASEDA 0251), cited at PHOIBLE: \url{https://phoible.org/sources/hercus_yarluyandi_nodate}
\item Wikipedia entry: \url{https://en.wikipedia.org/wiki/Yarluyandi}
\item AIATSIS Austlang record, which describes the Diyari word lists and lists Yarluyandi as closely related: \url{https://collection.aiatsis.gov.au/austlang/language/l17}
\end{itemize}

\section{Religious and Community Resources}
None found.

\section{Digital and Audio Resources}
The Mobile Language Team portal above is the only digital resource I found. I found no audio.

\section{Curriculum Development Context}
The language is dormant and documentation is limited to word lists. A heritage or revival curriculum would have to work with the Wangkangurru/Yarluyandi community and with the AIATSIS archive. It could also draw on better-documented relatives such as Diyari, subject to community wishes.
\end{document}